\begin{table}[H]
\centering
\caption{Top 10 Departamentos por Defección del Partido Nacional al Frente Amplio (Modelo con PI)}
\label{tab:dept_pn_top10}
\small
\begin{adjustbox}{max width=\textwidth}
\begin{tabular}{lrrrcc}
\toprule
\textbf{Departamento} & \textbf{N} & \textbf{Votos PN} & \textbf{PN$\rightarrow$FA (\%)} & \textbf{IC 95\%} & \textbf{Votos Est.} \\
\midrule
Artigas & 181 & 21,357 & 14.9 & [11.2, 18.6] & 3,185 \\
Cerro Largo & 208 & 28,909 & 14.7 & [12.4, 16.6] & 4,260 \\
Paysandú & 281 & 30,149 & 13.7 & [10.9, 16.3] & 4,118 \\
Treinta y Tres & 135 & 15,302 & 9.6 & [6.7, 12.4] & 1,465 \\
Soriano & 215 & 20,705 & 8.2 & [6.0, 10.7] & 1,700 \\
Flores & 59 & 8,333 & 6.6 & [1.4, 12.7] & 553 \\
Maldonado & 404 & 48,443 & 5.2 & [3.3, 7.2] & 2,538 \\
Tacuarembó & 234 & 21,940 & 3.9 & [1.8, 6.3] & 847 \\
Florida & 179 & 17,542 & 3.1 & [0.2, 6.5] & 547 \\
Lavalleja & 144 & 14,587 & 2.8 & [0.2, 6.0] & 407 \\
\bottomrule
\end{tabular}
\end{adjustbox}
\begin{flushleft}
\small \textit{Nota:} Estimaciones del modelo con Partido Independiente (PI) como categoría separada. PN, Partido Nacional; FA, Frente Amplio. Departamentos ordenados por probabilidad de transición PN$\rightarrow$FA. Votos estimados = votos PN primera vuelta $\times$ probabilidad de transición. Todos los modelos convergieron ($\hat{R} < 1.02$).
\end{flushleft}
\end{table}
